\chapter{Le Vocabulaire Fondamental du Risque}
\label{chap:vocrisque}

L'analyse de risque repose sur un vocabulaire précis. Une confusion entre ces termes conduit à des évaluations erronées.

\section{Le Triangle du Risque : La Modélisation}

Le risque est rarement isolé. Il résulte de la conjonction de trois facteurs fondamentaux, souvent représentés sous forme de triangle.

\begin{figure}[htbp]
\centering
\begin{tikzpicture}[
    node distance=2.5cm,
    every node/.style={align=center}
]
    % Nodes
    \node (asset) [draw=cyberblue, thick, fill=cyberblue!10, rounded corners, minimum width=3cm, minimum height=1cm] {\textbf{Actif} \\ (Ce que l'on protège)};
    \node (threat) [draw=cyberred, thick, fill=cyberred!10, rounded corners, minimum width=3cm, minimum height=1cm, below left=of asset] {\textbf{Menace} \\ (Ce qu'on craint)};
    \node (vuln) [draw=lizbizcolor, thick, fill=lizbizcolor!10, rounded corners, minimum width=3cm, minimum height=1cm, below right=of asset] {\textbf{Vulnérabilité} \\ (La faille)};

    % Arrows
    \draw[->, thick] (threat) -- node[above, sloped, font=\scriptsize] {Exploite} (vuln);
    \draw[->, thick] (vuln) -- node[right, font=\scriptsize] {Impacte} (asset);
    \draw[->, thick] (threat) -- node[left, font=\scriptsize] {Cible} (asset);
\end{tikzpicture}
\caption{Le triangle du risque}
\end{figure}

\subsection{L'Actif (Asset)}

\begin{boxdef}
\textbf{Actif :} Tout élément qui a de la valeur pour l'organisation et qui nécessite d'être protégé. Il peut être tangible (matériel) ou intangible (information, réputation).
\end{boxdef}

\textbf{Classification des actifs :}
\begin{itemize}
    \item \textbf{Primaires :} Activités vitales (Production, Vente).
    \item \textbf{Supports :} Aident à la production (RH, Comptabilité).
    \item \textbf{Infrastructure :} Fondations (Réseau, Électricité).
\end{itemize}

\begin{boxlizbiz}
\textbf{Chez LiZbiZ :}
L'ERP (logiciel de gestion) est un actif \textit{primaire} critique. S'il tombe, l'entrepôt s'arrête. La base de données clients est un actif \textit{intangible} sensible (RGPD).
\end{boxlizbiz}

\subsection{La Menace (Threat)}

\begin{boxdef}
\textbf{Menace :} Une cause potentielle d'incident, pouvant nuire à un actif. Elle est souvent caractérisée par son origine et sa motivation.
\end{boxdef}

\textbf{Sources de menace :}
\begin{enumerate}
    \item \textbf{Humaines :} Piratage externe, employé malveillant (interne), erreur humaine (négligence).
    \item \textbf{Environnementales :} Incendie, inondation, panne électrique.
\end{enumerate}

\begin{boxlizbiz}
\textbf{Chez LiZbiZ :} La principale menace est le \textit{Ransomware} (Rançonware). C'est une menace humaine externe, motivée par l'appât du gain financier.
\end{boxlizbiz}

\subsection{La Vulnérabilité (Vulnerability)}

\begin{boxdef}
\textbf{Vulnérabilité :} Une faiblesse ou un défaut dans un actif, un processus ou un contrôle, qui peut être exploité par une menace.
\end{boxdef}

\begin{boxalert}
Une menace sans vulnérabilité ne crée pas de risque (le feu n'est un risque que si le bâtiment est inflammable ou s'il n'y a pas d'extincteur).
\end{boxalert}

\begin{boxlizbiz}
\textbf{Chez LiZbiZ :} Le serveur Web est sous \textit{Windows Server 2012} non mis à jour. C'est une vulnérabilité technique critique (faille de sécurité connue).
\end{boxlizbiz}

% ------------------------------------------------------------------------------
\section{La Mesure du Risque : Probabilité et Gravité}
% ------------------------------------------------------------------------------

Pour évaluer le risque, on quantifie deux dimensions : la Probabilité et la Gravité.

\subsection{Probabilité (Likelihood)}

La \textbf{Probabilité} est la chance qu'un événement indésirable se produise.

\textbf{Échelle qualitative typique :}
\begin{enumerate}
    \item \textbf{Rare (1) :} Presque impossible.
    \item \textbf{Peu probable (2) :} Possible mais surprenant.
    \item \textbf{Possible (3) :} Peut se produire une fois par an.
    \item \textbf{Probable (4) :} Se produit régulièrement.
    \item \textbf{Quasi-certain (5) :} Se produit systématiquement.
\end{enumerate}

\begin{boxlizbiz}
\textbf{Chez LiZbiZ :} La probabilité d'une attaque par Ransomware sur le serveur Web obsolète est \textbf{Probable (4)}. C'est une cible facile sur Internet.
\end{boxlizbiz}

\subsection{Gravité ou Impact (Severity)}

La \textbf{Gravité} mesure l'ampleur des dommages causés à l'organisation si l'événement se réalise.

\textbf{Dimensions de l'impact :}
\begin{itemize}
    \item \textbf{Opérationnel :} Arrêt de la production, perte de temps.
    \item \textbf{Financier :} Perte d'argent directe, amendes.
    \item \textbf{Juridique :} Poursuites, non-conformité (RGPD).
    \item \textbf{Image :} Perte de confiance des clients, impact boursier.
\end{itemize}

\textbf{Échelle qualitative typique :}
\begin{enumerate}
    \item \textbf{Négligeable (1) :} Impact minime.
    \item \textbf{Limité (2) :} Gêne temporaire.
    \item \textbf{Significatif (3) :} Perte d'activité partielle.
    \item \textbf{Critique (4) :} Perte d'activité majeure.
    \item \textbf{Catastrophique (5) :} Mise en péril de l'entreprise.
\end{enumerate}

\begin{boxlizbiz}
\textbf{Chez LiZbiZ :} Si l'ERP tombe en panne (Ransomware), la gravité est \textbf{Catastrophique (5)}. Plus de commandes, plus de livraisons, perte de chiffre d'affaires immédiate et atteinte à l'image "livraison 24h".
\end{boxlizbiz}

\subsection{La Criticité (Risk Level)}

La \textbf{Criticité} (ou niveau de risque) est le produit de la Probabilité et de la Gravité.

\[ Criticité = Probabilité \times Gravité \]

On la visualise souvent via une \textbf{Matrice de Risque} (ou matrice de probabilité/gravité).

\begin{table}[htbp]
\centering
\begin{tabular}{|c|c|c|c|c|c|}
\hline
\textbf{Gravité / Prob.} & \textbf{1} & \textbf{2} & \textbf{3} & \textbf{4} & \textbf{5} \\ \hline
\textbf{5 (Catastroph.)} & Moyen & Élevé & Critique & \cellcolor{red!50}\textbf{Critique} & \cellcolor{red!50}\textbf{Critique} \\ \hline
\textbf{4 (Critique)} & Moyen & Élevé & Élevé & Critique & Critique \\ \hline
\textbf{3 (Signif.)} & Faible & Moyen & Élevé & Élevé & Critique \\ \hline
\textbf{2 (Limité)} & Faible & Faible & Moyen & Moyen & Élevé \\ \hline
\textbf{1 (Négligeable)} & Faible & Faible & Faible & Faible & Moyen \\ \hline
\end{tabular}
\caption{Matrice de criticité standard (1 à 5)}
\end{table}

% ------------------------------------------------------------------------------
\section{Risque Inhérent vs Risque Résiduel}
% ------------------------------------------------------------------------------

\subsection{Risque Inhérent (Inherent Risk)}

\begin{boxdef}
\textbf{Risque Inhérent :} Niveau de risque \textit{brut}, avant toute mise en place de mesures de protection (barrières techniques ou organisationnelles).
\end{boxdef}

C'est le risque "à nu". Il permet de juger de l'exposition naturelle de l'entreprise.

\subsection{Mesure de Sécurité (Safeguard / Countermeasure)}

Action technique (firewall, chiffrement) ou organisationnelle (charte, formation) mise en place pour réduire le risque.

\subsection{Risque Résiduel (Residual Risk)}

\begin{boxdef}
\textbf{Risque Résiduel :} Niveau de risque qui \textit{reste} après l'application des mesures de sécurité. C'est le risque accepté par l'entreprise.
\end{boxdef}

L'équation devient :
\[ Risque\_Résiduel = Risque\_Inhérent - Efficacité\_Mesures \]

\begin{boxlizbiz}
\textbf{Exemple LiZbiZ :}
\begin{itemize}
    \item \textbf{Risque Inhérent :} Panne électrique du serveur ERP (Probabilité 2, Gravité 5 $\rightarrow$ Risque Critique).
    \item \textbf{Mesure :} Installation d'un onduleur + Générateur électrique (Onglet "Business Continuity").
    \item \textbf{Risque Résiduel :} Panne ERP due à l'électricité (Probabilité 1, Gravité 1 $\rightarrow$ Risque Faible). Le risque est maîtrisé.
\end{itemize}
\end{boxlizbiz}


% ==============================================================================
% Fichier : 02_module_risque.tex (Extrait Chapitre 3)
% Description : Référentiels et Méthodes d'Analyse
% ==============================================================================